\begin{center}
\begin{minipage}{0.90\textwidth}
\small
\textbf{Abstract}
\vspace{3pt}

Vision-language models (VLMs) built by concatenating visual tokens into a
frozen or fine-tuned LLM pay a steep price: attention cost grows with image
resolution, and adapting the backbone erodes the language model's original
text reasoning. Parameter-rotation methods remove visual tokens from the
backbone, but still depend on a heavy external vision tower, and they
deteriorate once the model generates autoregressively on its own. We introduce
\modelname, a framework that converts a completely frozen pure-text LLM into a
VLM that keeps \textbf{zero visual tokens in the backbone} and, in its
strongest configuration, \textbf{removes the external vision tower entirely}.
\modelname{} rests on two mechanisms. A \emph{dual-tier orthogonal rotation}
architecture first repurposes the student's own first $K{=}4$ layers as a
built-in vision encoder through static Cayley rotation and channel-wise
scaling ($2.4$M parameters), then discards every image token and modulates the
deeper frozen layers with dynamic block-diagonal rotations generated by
multi-frequency or spatial-anchor queries. Because the text pathway is left
untouched, pure-text behaviour is bit-exact, whereas merging the same vision
adapters into the backbone collapses GSM8K from $70.4\%$ to $0.0\%$.
Second, \emph{online policy distillation} (VL-OPD) supervises the student's
own rollouts with token-level KL from a stronger VLM, which lifts free-form
generation by up to $+9.9$ points. With the same training data, steps, and
frozen base model, \modelname{} outperforms a token-concatenation LLaVA
baseline on visual perception and anti-hallucination
($44.37\%$ vs.\ $41.35\%$), matches it on the average of perception,
hallucination and mathematical reasoning ($35.90\%$ vs.\ $35.87\%$), is
strictly better than single-frequency rotation and LoRA-based built-in
encoders, and does so with $74.2\%$ fewer FLOPs and $100\%$ preserved text
ability.
\end{minipage}
\end{center}
